\documentclass[11pt,a4paper]{article}
\usepackage[utf8]{inputenc}
\usepackage[french]{babel}
\usepackage[T1]{fontenc}
\usepackage{amsmath}
\usepackage{amsfonts}
\usepackage{amssymb}
\usepackage{graphicx}
\usepackage{pdfpages}
\usepackage{palatino}
\usepackage{wrapfig} 
\usepackage{xcolor}
\definecolor{lightgrey}{gray}{0.8}
\usepackage{listingsutf8}
\lstset
{ language=java
, captionpos=b
%, frame=lines
, morekeywords={var,get,set}
, basicstyle=\footnotesize\ttfamily
, keywordstyle=\color{blue}
, commentstyle=\color{green}
, stringstyle=\color{red}
, numbers=left
, numberstyle=\scriptsize
, stepnumber=2
, numbersep=5pt
, breaklines=true
, tabsize=2
, showstringspaces=false
, emph={double,bool,int,unsigned,char,true,false,void,get,set}
, emphstyle=\color{blue}
, emph={Assert,Test}
, emphstyle=\color{red}
, emph={[2]\#using,\#define,\#ifdef,\#endif}
, emphstyle={[2]\color{blue}}
, frame=single
, rulesepcolor=\color{lightgrey}
, lineskip={-1.5pt} % single line spacing
}
\usepackage[left=2cm,right=2cm,top=2cm,bottom=2cm]{geometry}
\title{SINF1121 - Groupe 9\\Rapport 1}
\author{\\Gégo Anthony\\ Gena Xavier \\Joveneau Quentin\\Libioulle Thibault \\Moyaux Arnold\\ Naveau Adrien \\ Payen Marlon}
\def\blurb{\textsc{Université catholique de Louvain\\
  École polytechnique de Louvain}}
\def\clap#1{\hbox to 0pt{\hss #1\hss}}%
\def\ligne#1{%
  \hbox to \hsize{%
    \vbox{\centering #1}}}%
\def\haut#1#2#3{%
  \hbox to \hsize{%
    \rlap{\vtop{\raggedright #1}}%
    \hss
    \clap{\vbox{\vfill\centering #2\vfill}}%
    \hss
    \llap{\vtop{\raggedleft #3}}}}%
\begin{document}
\begin{titlepage}
\thispagestyle{empty}\vbox to 1\vsize{%
  \vss
  \vbox to 1\vsize{%
    \haut{\includegraphics[scale=0.15]{logo_ucl.pdf}}{\blurb}{\includegraphics[scale=0.4]{logo_epl.jpg}}
    \vfill
    \ligne{\Huge \textbf{\textsc{Structures de données\\ et algorithmes (SINF1121)}}}
    \vspace{5mm}
    \ligne{\Large \textbf{Mission 1 - Réponses aux questions}}
    \vspace{5mm}
    \ligne{\large{-- 25 septembre 2013 --}}
    %\begin{center}\includegraphics[scale=3]{img/img_couverture.png}\end{center}
    \vfill
    \ligne{%
      \begin{tabular}{c}
        \textsc{Travail du groupe 9 :}
      \end{tabular}}
    \vspace{5mm}
    \ligne{%
      \begin{tabular}{lrclr}
         \textsc{Gégo} Anthony  & 28581100 & \hspace{80pt} & \textsc{Moyau} Arnold & xxxxxx00\\
         \textsc{Gena} Xavier  & xxxxxx00 & & \textsc{Naveau} Adrien & xxxxxx00\\
         \textsc{Joveneau} Quentin & xxxxxx00  & & \textsc{Payen} Marlon & xxxxxx00\\
         \textsc{Libioulle} Thibault & xxxxxx00  & &  & 
      \end{tabular}
      }
    }%
  \vss
  }
\end{titlepage}

\section{Question 1}
\begin{quotation}
\textit{Définissez ce qu’est un type abstrait de données (TAD). En Java, est-il préférable de décrire un TAD par une classe ou une interface ? Pourquoi ?}
\end{quotation}

Un type abstrait de données est un modèle mathématique d'une structure de données spécifiant le type d'informations stockées, les opérations supportées sur ces dernières, ainsi que les types de paramètres des opérations. Un TAD spécifie ce que chaque opération fait, mais pas comment elle le fait. 

De ce fait, en Java, un TAD peut être décrit au moyen d'une interface, puisqu'elle ne consiste qu'en une liste de déclarations de méthodes sans implémentation.

\section{Question 2} 
\begin{quotation}
\textit{Comment faire pour implémenter une pile par une liste simplement chaînée où les
opérations push et pop se font en fin de liste ? Cette solution est-elle efficace ?
Argumentez.}
\end{quotation}

L'essentiel est de maintenir une référence vers l'objet en fin de liste. Chaque nœud possède alors une référence vers l'élément contenu ainsi que vers le nœud qui le précède.

L'opération push se réalisera, quelque soit la taille de la pile, en un nombre constant d'opérations. Il suffit en effet que de lié le nouveau nœud à l'actuel dernier, et de modifier la référence pointant vers le dernier nœud vers le nouveau. Un raisonnement similaire peut être appliqué pour l'opération pop qui consiste à enlever le dernier nœud et modifier la référence du dernier nœud vers le précédent.

Au vu de ce raisonnement, nous pouvons conclure que la solution sera tout aussi efficace que si l'opération était effectuée en début de liste.

\section{Question 3}
\begin{quotation}
\textit{En consultant la documentation sur l’API de Java, décrivez l’implémentation
d’une pile par la classe java.util.Stack. Cette classe peut-elle convenir
comme implémentation de l’interface Stack décrite dans DSAJ-5 ? Pourquoi ?}
\end{quotation}

La classe java.util.Stack étend la classe Vector à laquelle sont ajoutées les opérations nécessaires pour être utilisable comme pile. Un Vector est un tableau qui est redimensionné dynamiquement en fonction des besoins. Si, d'un point de vue, la complexité spatiale peut être réduite car il n'est pas nécessaire de retenir les références vers les noeuds suivants, il ne faut cependant pas négliger les opérations nécessaire pour la réallocation du tableau à chaque redimensionnement. 

Puisque l'interface est complètement indépendante de l'implémentation et que les méthodes push et pop sont définies, la classe peut ainsi servir d'implémentation de l'interface Stack décrite dans le livre de référence.

\section{Question 4}
\begin{quotation}
\textit{Proposez une implémentation de la classe DNodeStack. Il s’agit d’une classe
similaire à NodeStack (décrite dans DSAJ-5) qui propose une implémentation
générique d’une pile. Votre classe DNodeStack doit utiliser une implémentation en liste doublement chaînée générique. Elle reposera donc sur une classe
DNode<E> (similaire à Node<E>, décrite dans DSAJ-5) que vous devez définir.
Ajoutez, dans la classe DNodeStack, une méthode public String toString()
qui renvoie une chaîne de caractères représentant le contenu de la pile. Commentez votre code.}
\end{quotation}

\lstinputlisting[caption=DNode.java]{question4/DNode.java}
\lstinputlisting[caption=Stack.java]{question4/Stack.java}
\lstinputlisting[caption=DNodeStack.java]{question4/DNodeStack.java}
\lstinputlisting[caption=EmptyStackException.java]{question4/EmptyStackException.java}
\end{document}